% !TEX root = chapter-standalone.tex

\section{Bicategories}
\label{sec:bicats-def}

% \linkvideo{...} % TODO: add video link

Throughout this text, we have encountered many situations in which we have objects, morphisms between objects, and composition of morphisms --- that is, categories.
But we have also seen, in more than one place, that the story does not end with morphisms: there can also be meaningful \emph{transformations between morphisms}.
We met this idea when we studied \SY{natural transformations} between \SY{functors}: given two functors~$\funa, \funb \colon \CatC \fto \CatD$, there may exist a natural transformation~$\ntrafoa \colon \funa \nto \funb$.


The notion of a \maindef{bicategory} is designed to capture exactly this kind of structure in its full generality.
A bicategory is a setting in which we have:
\begin{itemize}
    \item objects (which we also call \emph{0-cells});
    \item morphisms between objects (called \emph{1-cells}); and
    \item transformations between morphisms (called \emph{2-cells}).
\end{itemize}

The key novelty, compared to an ordinary category, is the presence of the 2-cells and the fact that there are now \emph{two} kinds of composition for them: one that composes 2-cells ``vertically'' (one after the other, when they share a common 1-cell boundary) and one that composes them ``horizontally'' (side by side, when the 1-cells they relate are themselves composable).
Furthermore, in a bicategory, the composition of 1-cells is only required to be associative and unital \emph{up to coherent isomorphism}, rather than strictly.
This weakening is crucial: most of the naturally occurring examples of ``categories with an extra layer'' have this weaker structure.

\subsection{Motivation: why bicategories?}

Before giving the formal definition, let us discuss why one might want such a concept in the first place.

There is a simple pattern that recurs in many areas of mathematics and its applications: one has a collection of ``spaces'' of some kind, and between any two of them, there is not just a set of maps, but a \emph{category} of maps --- that is, the maps themselves admit transformations between them.

\begin{example}[Categories, functors, and natural transformations]
    \label{exa:bicat-motivation-cat}
    Consider the world of (small) \SY{categories}.
    The objects are categories, the 1-cells are \SY{functors}, and the 2-cells are \SY{natural transformations}.
    We know that \SY{functors} can be composed (this was defined in the section on functors), and we also know that \SY{natural transformations} can be composed, both ``vertically'' (yielding a natural transformation between the composite functors) and ``horizontally'' (by whiskering or the Godement product).
    However, one must be careful: functor composition is strictly associative and unital, so this particular example is actually a \emph{strict} 2-category (meaning a bicategory in which the coherence isomorphisms are all identities).
\end{example}

\begin{example}[Profunctors]
    \label{exa:bicat-motivation-prof}
    Perhaps the most important example for us in this book is the world of \SY{profunctors}.
    We have already defined profunctor composition (see \cref{def:profunctor-composition}) and noted that it is \emph{not} strictly associative --- it is only associative up to a natural isomorphism.
    Hence profunctors do not form an ordinary category.
    This is precisely the situation that bicategories are designed to handle.
    The bicategory of profunctors has: categories as objects, profunctors as 1-cells, and natural transformations between profunctors as 2-cells.
\end{example}

\begin{example}[Relations]
    \label{exa:bicat-motivation-rel}
    Consider sets as objects and relations~$R \subseteq \setA \cartprod \setB$ as 1-cells from~$\setA$ to~$\setB$.
    Between two relations~$R, S \subseteq \setA \cartprod \setB$, we can consider the inclusion~$R \subseteq S$ as a 2-cell.
    Composition of relations (which we have studied earlier) makes this into a bicategory.
    In fact, since there is at most one 2-cell between any two parallel 1-cells (inclusion is a property, not additional data), this is a \emph{locally posetal} bicategory.
\end{example}

These examples show that bicategories arise naturally whenever one ``enriches'' the morphisms of a category with extra structure.
The formal definition below axiomatizes the essential features shared by all of these examples.

\subsection{The definition}

We now state the full definition. It involves a fair amount of data. We recommend reading through it once to get the overall shape, and then revisiting the examples that follow.

\begin{ctdefinition}[Bicategory]
    \SYNDEF{bicategory}
    \label{def:bicategory}
    A \maindef{bicategory}~$\CatB$ consists of:

    \constit

    \begin{enumerate}
        \item \label{bicat-objects}
              \emph{Objects (0-cells):}
              A collection~$\Ob(\CatB)$ of objects.

        \item \label{bicat-hom-cats}
              \emph{Hom-categories (1-cells and 2-cells):}
              For every pair of objects~$\Obja, \Objb \setin \Ob(\CatB)$, a category~$\CatB(\Obja, \Objb)$, called a \emph{hom-category}.
              The objects of~$\CatB(\Obja, \Objb)$ are called \emph{1-cells} (or \emph{morphisms}) and are denoted~$\mora \colon \Obja \mto \Objb$.
              The morphisms of~$\CatB(\Obja, \Objb)$ are called \emph{2-cells} and are denoted~$\ntrafoa \colon \mora \nto \morb$ for~$\mora, \morb \colon \Obja \mto \Objb$.
              Composition of 2-cells within a hom-category is called \emph{vertical composition} and is denoted~$\ntrafoa \ntrafothen \ntrafob$ for composable 2-cells~$\ntrafoa$ and~$\ntrafob$.
              The identity 2-cell on a 1-cell~$\mora$ is denoted~$\catidof{\mora}$.

        \item \label{bicat-horizontal-comp}
              \emph{Horizontal composition:}
              For every triple of objects~$\Obja, \Objb, \Objc \setin \Ob(\CatB)$, a functor
              \begin{equation}\label{eq:bicat-horiz-comp}
                  \mthenof{\Obja,\Objb,\Objc} \colon \CatB(\Obja, \Objb) \Ctimes \CatB(\Objb, \Objc) \fto \CatB(\Obja, \Objc),
              \end{equation}
              called \emph{horizontal composition}.
              We write~$\mora \mthen \morb$ for the horizontal composite of 1-cells~$\mora\colon \Obja \mto \Objb$ and~$\morb\colon \Objb \mto \Objc$, and~$\ntrafoa \nthenhor \ntrafob$ for the horizontal composite of 2-cells.

        \item \label{bicat-identities}
              \emph{Identity 1-cells:}
              For every object~$\Obja \setin \Ob(\CatB)$, a distinguished 1-cell
              \begin{equation}\label{eq:bicat-identity-cell}
                  \catidof{\Obja} \colon \Obja \mto \Obja,
              \end{equation}
              called the \emph{identity} 1-cell at~$\Obja$.

        \item \label{bicat-associator}
              \emph{Associator:}
              For every composable triple of 1-cells~$\mora\colon \Obja \mto \Objb$, $\morb\colon \Objb \mto \Objc$, $\morc\colon \Objc \mto \Objd$, an invertible 2-cell
              \begin{equation}\label{eq:bicat-associator}
                  \associator_{\mora,\morb,\morc} \colon (\mora \mthen \morb) \mthen \morc \nto \mora \mthen (\morb \mthen \morc),
              \end{equation}
              called the \emph{associator}.

        \item \label{bicat-unitors}
              \emph{Left and right unitors:}
              For every 1-cell~$\mora \colon \Obja \mto \Objb$, invertible 2-cells
              \begin{equation}\label{eq:bicat-left-unitor}
                  \leftunitor_\mora \colon \catidof{\Obja} \mthen \mora \nto \mora,
              \end{equation}
              \begin{equation}\label{eq:bicat-right-unitor}
                  \rightunitor_\mora \colon \mora \mthen \catidof{\Objb} \nto \mora,
              \end{equation}
              called the \emph{left unitor} and \emph{right unitor}, respectively.
    \end{enumerate}

    \condit

    \begin{enumerate}
        \item \label{bicat-naturality}
              \emph{Naturality:}
              The associator~$\associator$ and the unitors~$\leftunitor$,~$\rightunitor$ are \SY{natural} in their arguments.
              (That is, they are natural isomorphisms, not merely families of isomorphisms.)

        \item \label{bicat-pentagon}
              \emph{Pentagon axiom:}
              For every composable quadruple of 1-cells~$\mora, \morb, \morc, \mord$, the following diagram of 2-cells commutes:
              % TODO: \equationsag{bicat_pentagon}{eq:bicat-pentagon}
              \begin{equation}\label{eq:bicat-pentagon}
                  \begin{tikzcd}[column sep=small, row sep=large]
                      & {((\mora \mthen \morb) \mthen \morc) \mthen \mord}
                      \ar[dl, "{\associator_{\mora,\morb,\morc}\, \nthenhor\, \catidof{\mord}}"']
                      \ar[dr, "{\associator_{\mora \mthen \morb,\morc,\mord}}"]
                      & \\
                      {(\mora \mthen (\morb \mthen \morc)) \mthen \mord}
                      \ar[d, "{\associator_{\mora,\morb \mthen \morc,\mord}}"']
                      & &
                      {(\mora \mthen \morb) \mthen (\morc \mthen \mord)}
                      \ar[d, "{\associator_{\mora,\morb,\morc \mthen \mord}}"] \\
                      {\mora \mthen ((\morb \mthen \morc) \mthen \mord)}
                      \ar[rr, "{\catidof{\mora}\, \nthenhor\, \associator_{\morb,\morc,\mord}}"']
                      & &
                      {\mora \mthen (\morb \mthen (\morc \mthen \mord))}
                  \end{tikzcd}
              \end{equation}

        \item \label{bicat-triangle}
              \emph{Triangle axiom:}
              For every composable pair of 1-cells~$\mora \colon \Obja \mto \Objb$ and~$\morb \colon \Objb \mto \Objc$, the following diagram of 2-cells commutes:
              % TODO: \equationsag{bicat_triangle}{eq:bicat-triangle}
              \begin{equation}\label{eq:bicat-triangle}
                  \begin{tikzcd}[column sep=large]
                      {(\mora \mthen \catidof{\Objb}) \mthen \morb}
                      \ar[rr, "{\associator_{\mora, \catidof{\Objb}, \morb}}"]
                      \ar[dr, "{\rightunitor_\mora\, \nthenhor\, \catidof{\morb}}"']
                      & &
                      {\mora \mthen (\catidof{\Objb} \mthen \morb)}
                      \ar[dl, "{\catidof{\mora}\, \nthenhor\, \leftunitor_\morb}"] \\
                      & {\mora \mthen \morb} &
                  \end{tikzcd}
              \end{equation}
    \end{enumerate}

\end{ctdefinition}

\begin{remark}
    Notice that there is a strong resemblance with the definition of \SY{monoidal categories}: a monoidal category also has an associator and unitors satisfying pentagon and triangle axioms.
    This is no coincidence.
    A \SY{monoidal category} is precisely a bicategory with exactly one object.
    Just as a \SY{monoid} is a category with one object, a \SY{monoidal category} is a bicategory with one object.
    This is an important conceptual point: it means that the theory of monoidal categories is subsumed by the theory of bicategories.
\end{remark}

\begin{remark}
    \label{rem:bicat-why-weak}
    The reason we require associativity and unitality only up to isomorphism (rather than on the nose) is that this is what we encounter in practice.
    As we noted for \SY{profunctors}, their composition is defined via a colimit (a quotient of a coproduct), and colimits are typically only unique up to isomorphism, not up to equality.
    More generally, in almost all ``naturally occurring'' higher-categorical structures, the algebraic operations are only well-defined up to a coherent system of isomorphisms.
    The pentagon and triangle axioms ensure that these isomorphisms fit together consistently, so that one can unambiguously compose any finite sequence of 1-cells.
    (This is the content of Mac~Lane's coherence theorem for monoidal categories, which extends to bicategories.)
\end{remark}

\subsection{Vertical and horizontal composition}
\label{sec:bicat-compositions}

Let us unpack the two kinds of composition that appear in a bicategory.

\emph{Vertical composition.}
Given 2-cells~$\ntrafoa \colon \mora \nto \morb$ and~$\ntrafob \colon \morb \nto \morc$ in a hom-category~$\CatB(\Obja, \Objb)$, their \emph{vertical composite} is the 2-cell
\begin{equation}
    \ntrafoa \ntrafothen \ntrafob \colon \mora \nto \morc.
\end{equation}
This is just ordinary composition in the category~$\CatB(\Obja, \Objb)$.
In particular, it is strictly associative and unital (the identity 2-cells~$\catidof{\mora}$ serve as identities).

\emph{Horizontal composition.}
Given 2-cells~$\ntrafoa \colon \mora \nto \mora'$ in~$\CatB(\Obja, \Objb)$ and~$\ntrafob \colon \morb \nto \morb'$ in~$\CatB(\Objb, \Objc)$, their \emph{horizontal composite} is the 2-cell
\begin{equation}
    \ntrafoa \nthenhor \ntrafob \colon \mora \mthen \morb \nto \mora' \mthen \morb'.
\end{equation}
This is well-defined because horizontal composition~$\mthenof{\Obja,\Objb,\Objc}$ is required to be a \emph{functor}, and a functor acts on both objects and morphisms.
In particular, functoriality implies the \emph{interchange law}: for composable 2-cells~$\ntrafoa, \ntrafoa'$ (vertical) and~$\ntrafob, \ntrafob'$ (vertical), we have
\begin{equation}\label{eq:bicat-interchange}
    (\ntrafoa \ntrafothen \ntrafoa') \nthenhor (\ntrafob \ntrafothen \ntrafob') = (\ntrafoa \nthenhor \ntrafob) \ntrafothen (\ntrafoa' \nthenhor \ntrafob').
\end{equation}
This says that when we have a grid of 2-cells, the result of composing them does not depend on whether we compose vertically first and then horizontally, or the other way around.

\subsection{Strict 2-categories}
\label{sec:strict-2cat}

An important special case deserves its own name.

\begin{ctdefinition}[Strict 2-category]
    \label{def:strict-2-category}
    A \maindef{strict 2-category} is a bicategory in which the associator and unitors are all identity 2-cells.
    Equivalently, it is a bicategory in which horizontal composition of 1-cells is strictly associative and strictly unital:
    \begin{equation}\label{eq:strict-2cat-assoc}
        (\mora \mthen \morb) \mthen \morc = \mora \mthen (\morb \mthen \morc), \qquad \catidof{\Obja} \mthen \mora = \mora = \mora \mthen \catidof{\Objb}.
    \end{equation}
\end{ctdefinition}

Every strict 2-category is, in particular, a bicategory (with all coherence isomorphisms being identities).
But not every bicategory is equivalent to a strict 2-category --- the weakness is essential.

\subsection{Examples}
\label{sec:bicat-examples}

We now revisit and expand the motivating examples from the beginning of this section.

\begin{example}[\Category as a strict 2-category]
    \label{exa:bicat-cat}
    The strict 2-category~$\Category$ has:
    \begin{itemize}
        \item \emph{Objects:} small categories;
        \item \emph{1-cells:} \SY{functors} between categories;
        \item \emph{2-cells:} \SY{natural transformations} between functors.
    \end{itemize}
    Vertical composition of 2-cells is the usual (componentwise) composition of natural transformations.
    Horizontal composition is given by the Godement product (horizontal composition of natural transformations).
    Functor composition is strictly associative and unital, so the associators and unitors are identities, making this a strict 2-category.
\end{example}

\begin{example}[The bicategory of profunctors]
    \label{exa:bicat-prof}
    The bicategory~$\Prof$ has:
    \begin{itemize}
        \item \emph{Objects:} small categories;
        \item \emph{1-cells:} \SY{profunctors}~$\funa \colon \CatC \profto \CatD$;
        \item \emph{2-cells:} natural transformations between profunctors (viewed as functors~$\CatCop \Ctimes \CatD \fto \Set$).
    \end{itemize}
    Horizontal composition of 1-cells is profunctor composition (see~\cref{def:profunctor-composition}), which involves a coend (colimit) construction.
    The identity 1-cell on a category~$\CatC$ is the hom-functor~$\Hom_\CatC \colon \CatCop \Ctimes \CatC \fto \Set$.
    Because the coend is only defined up to isomorphism, the associators and unitors are nontrivial --- this is a genuine bicategory, not a strict 2-category.
\end{example}

\begin{example}[The bicategory of spans]
    \label{exa:bicat-span}
    Fix a category~$\CatC$ with pullbacks (for example, $\CatC = \Set$).
    The bicategory~$\operatorname{Span}(\CatC)$ has:
    \begin{itemize}
        \item \emph{Objects:} objects of~$\CatC$;
        \item \emph{1-cells from~$\Obja$ to~$\Objb$:} spans, i.e., diagrams~$\Obja \leftarrow \Objc \rightarrow \Objb$ in~$\CatC$;
        \item \emph{2-cells:} morphisms of spans, i.e., morphisms~$\Objc \mto \Objc'$ making the obvious triangles commute.
    \end{itemize}
    Composition of spans is given by pullback, which is only defined up to isomorphism.
    Hence the associator is nontrivial, and this is a genuine bicategory.
\end{example}

\begin{example}[Monoidal categories as one-object bicategories]
    \label{exa:bicat-monoidal}
    Let~$\tup{\CatC, \mtimescat, \idmoncat}$ be a \SY{monoidal category}.
    We can construct a bicategory~$\mathbf{\Sigma}\CatC$ with:
    \begin{itemize}
        \item a single object~$\star$;
        \item 1-cells: objects of~$\CatC$ (so $\mathbf{\Sigma}\CatC(\star, \star) = \CatC$);
        \item 2-cells: morphisms of~$\CatC$.
    \end{itemize}
    Horizontal composition of 1-cells is the monoidal product~$\mtimescat$.
    The identity 1-cell is the monoidal unit~$\idmoncat$.
    The associator and unitors of the bicategory are precisely the associator and unitors of the monoidal category.
    Conversely, given a bicategory with a single object, the hom-category at that object inherits a monoidal structure.
    This establishes the precise correspondence: \emph{monoidal categories are one-object bicategories}.
\end{example}


\begin{remark}
    Among the examples above, the bicategory of profunctors is the most important for the applications discussed in this book.
    The fact that profunctors do not form a category, but do form a bicategory, was noted at the end of the section on profunctor composition.
    Bicategories give us the right framework to study profunctors and their composition coherently.
\end{remark}

\subsection{Looking ahead: why bicategories matter for applications}
\label{sec:bicat-motivation-applications}

We conclude this section by briefly indicating why bicategories are useful beyond pure mathematics, especially in the context of the engineering applications discussed in this book.

Firstly, many of the structures we have studied --- \SY{design problems}, \SY{profunctors}, relations, and spans --- naturally live in bicategories.
The bicategorical perspective gives us a uniform language for discussing composition, identity, and coherence in all of these settings.

Secondly, many important categorical concepts generalize from categories to bicategories.
In particular, the notions of \emph{adjunction} and \emph{monad} make sense in any bicategory, and their bicategorical versions unify several constructions that might otherwise seem unrelated.
We will explore this in subsequent sections.

Thirdly, the relationship between monoidal categories and one-object bicategories (\cref{exa:bicat-monoidal}) means that any theorem about bicategories automatically gives a theorem about monoidal categories.
Since monoidal categories are central to many applications (modeling parallel composition, resource theories, type systems, and more), this is a powerful source of general results.